\chapter{Conclusion}\label{ch:conclusion}

\section{Conclusion}\label{sec:conclusion}

Across \NumCells{} pre-registered clinical cells \rev{and $12$ external validation cells},
\shap{} and \lime{} agree substantially ($\rho=\MedianRho{}$,
\rev{$p_{\text{adj}}<10^{-3}$}) without being interchangeable, \rev{with disagreement
heightened by feature redundancy rather than task difficulty.} On tree models,
\textsc{TreeSHAP} dominates speed \rev{across all \lime{} budgets ($500$--$10{,}000$)},
local robustness, and local faithfulness---\rev{including conditional $k$-NN deletion
(\CondShapWins{}/18 primary, \ScaleCondWins{}/12 external, $p_{\text{adj}}<0.05$)---while
\lime{} is faster on stacking and ties \shap{} under global \textsc{roar} retraining
($p=\RoarAucP{}$).} Stability must be reported across all three axes because retraining
shifts deterministic \textsc{TreeSHAP} ($\rho=\RetrainShap{}$), and rank fusion never beats
the best single explainer (\FusionMaxWins{}/18 primary, \rev{\ScaleFusWins{}/12 external})
given minimal oracle headroom.


\section{Future work}
Four extensions follow directly. First, replacing the literature-derived plausibility
reference with synthetic data whose true importances are known by construction, which
converts \RQ{5} from a proxy into a ground-truth check. Second, extending the
comparison to counterfactual and rule-based explainers, which answer a different
question and may prove complementary in the way SHAP and LIME did not. Third, an
on-manifold fidelity evaluation using a generative model of the data, removing the
independence assumption both of our baselines share. Fourth, a human study: every metric
here is a property of the model--explanation pair, and whether either explanation helps
a person decide better is untested.
